# Frozen main task: H^s operator-domain vs abstract completion

Let `K_c = -d^2/dx^2 + c` on `[-1,1]` with Krein boundary condition
`f'(±1) = (f(1)-f(-1))/2`, `c > 0`. Let `H^s`, `s >= 4`, be the left-definite
space associated with `K_c`. Let `{Q_n^(s)}` be the SL_hs orthogonal polynomial
system defined via the isometries `K_c^{-r}` on `L^2` or `H^1`.

Prove or disprove, for integer `s >= 4`:

1. Give a necessary and sufficient condition for `Q_n^(s) ∈ D(K_c^(s/2))`.
2. Determine whether the operator-domain completion `D(K_c^(s/2))` equals the
   abstract completion obtained from the left-definite inner product on
   polynomials.
3. Determine whether `span{Q_n^(s)}` is dense in `D(K_c^(s/2))` under the
   operator-domain reading.

The complete polynomial degree spectrum is a bonus, not a completion gate.

Rules: do not inspect repository history, current project files, known solution,
or network. State all external theorems with hypotheses. Numerical evidence is
not proof.
